\documentclass{article}
\usepackage{amsmath, amssymb}
\newcommand{\PW}{P(W^{1/2})}

\title{Geodesic Newton Direction with Blended RHS}
\author{}
\date{}

\begin{document}
\maketitle

\section{Setup}

Let $b_0$ denote the original affine term, $c$ the original cost, $e$ the
identity element, $A$ the constraint operator, and $P(X)$ the quadratic
representation $P(X)v = X v X$.  The Gram matrix is
\[
  G = A^T P(W) A.
\]

\section{Three-term decomposition}

The Newton step at the current scaling $W$ is decomposed into three
components that are independent of any blending parameters.

\subsection{Gram solves}

\[
  G\,y_0 = A^T(2W),
\]
\[
  G\,y_1^{(0)} = -\bigl(c + A^T P(W)\,b_0\bigr),
\]
\[
  G\,y_1^{(\theta)} = c + A^T P(W)\,b_0 - A^T\bigl(e + P(W)\,e\bigr).
\]

\subsection{Direction components}

\[
  d_0 = e - \PW(A\,y_0),
\]
\[
  d_1^{(0)} = \PW(-b_0 - A\,y_1^{(0)}),
\]
\[
  d_1^{(\theta)} = \PW(b_0 - e - A\,y_1^{(\theta)}).
\]

All six quantities ($y_0, y_1^{(0)}, y_1^{(\theta)}, d_0, d_1^{(0)},
d_1^{(\theta)}$) depend only on $W$, the original data $(b_0, c)$, and
the constraint operator $A$.  They are computed with one factorization
and three back-solves.

\section{Parameterized direction}

Given independent parameters $\tau$ (original-problem weight) and
$\theta$ (identity-centering weight), the Newton direction is
\[
  \boxed{d(k,\tau,\theta) = d_0 + k\bigl(\tau\,d_1^{(0)}
    + \theta\,d_1^{(\theta)}\bigr),}
\]
which is affine in $k$, $\tau$, and $\theta$.

The corresponding dual variable is
\[
  y(k,\tau,\theta) = y_0 + k\bigl(\tau\,y_1^{(0)}
    + \theta\,y_1^{(\theta)}\bigr).
\]

\subsection{Special cases}

\begin{itemize}
  \item $\tau = 1,\;\theta = 0$: original problem,
        $d = d_0 + k\,d_1^{(0)}$.
  \item $\tau = 0,\;\theta = 1$: feasibility/centering problem with
        $b = e$, $c = A^T e$, \; $d = d_0 + k\,d_1^{(\theta)}$.
  \item $\tau = 1-\theta$: single-parameter blend
        $b_\theta = \theta\,e + (1-\theta)\,b_0$, \;
        $c_\theta = \theta\,A^T e + (1-\theta)\,c$.
\end{itemize}

\section{Minimum-norm $k$ at fixed $(\tau,\theta)$}

Define $\tilde{d}_1 = \tau\,d_1^{(0)} + \theta\,d_1^{(\theta)}$.
The squared norm of the direction is
\[
  \|d(k)\|^2 = \|d_0\|^2 + 2k\,\langle d_0,\,\tilde{d}_1\rangle
               + k^2\,\|\tilde{d}_1\|^2,
\]
which is minimized at
\[
  \boxed{k^*(\tau,\theta)
    = -\frac{\langle d_0,\, \tau\,d_1^{(0)}
             + \theta\,d_1^{(\theta)}\rangle}
            {\|\tau\,d_1^{(0)}
             + \theta\,d_1^{(\theta)}\|^2}.}
\]

\section{Slack, primal variable, and multiplier}

The Newton direction $d$ lives in $v$-space, where $W = \exp(v)$.
The linearized primal slack and dual multiplier are recovered from $d$
as follows.

\subsection{Slack}

The primal slack (cone element) is
\[
  \boxed{s = \frac{1}{k}\,P(W^{-1/2})(e - d).}
\]

\subsection{Dual multiplier}

The dual multiplier is
\[
  \boxed{\lambda = \frac{1}{k}\,P(W^{1/2})(e + d)
    = \frac{1}{k}\,W\circ(e + d).}
\]

At the central path ($d = 0$), $s = \lambda = W/k$ and
$s \circ \lambda = W^2/k^2 = \mu\,W^2$, where $\mu = 1/k^2$.

\subsection{Primal variable}

The optimization variable is
\[
  \boxed{x = \frac{y}{k}
    = \frac{1}{k}\,y_0 + \tau\,y_1^{(0)} + \theta\,y_1^{(\theta)}.}
\]

\subsection{Complementarity}

The complementarity gap is
\[
  \langle s, \lambda \rangle
  = \frac{1}{k^2}\,\langle P(W^{1/2})(e-d),\, P(W^{1/2})(e+d)\rangle
  = \mu\,\bigl(m - \|d\|^2\bigr),
\]
where $m$ is the rank (dimension of the cone) and
$\mu = 1/k^2$.  At the central path ($d=0$), the gap is $\mu\,m$.

\section{Duality residual}

The Newton direction satisfies the identity
\[
  \boxed{b_0^T \lambda + c^T x + \frac{\mu}{\tau}
    = \theta\,\bigl(\langle b_0, e\rangle + 1\bigr),}
\]
where $\lambda$, $x$, $\mu = 1/k^2$ are as defined above.
This can be verified by substituting the expressions for $\lambda$
and $x$ in terms of the decomposition.

At the central path ($d = 0$, $\tau = 1$, $\theta = 0$), the left-hand
side reduces to $b_0^T \lambda + c^T x + \mu$, and the right-hand
side vanishes, recovering the standard complementarity relation.

At the trivially centered point ($\theta = 1$, $\tau = 1$, $k = 1$,
$W = e$, $d = 0$): $\lambda = e$, $x = 0$, $\mu = 1$, and
\[
  b_0^T e + 0 + 1 = \langle b_0, e\rangle + 1,
\]
which is satisfied identically.

The \emph{duality residual} is the absolute error in this equation:
\[
  r_{\mathrm{dual}} = \bigl|b_0^T \lambda + c^T x + \tfrac{\mu}{\tau}
    - \theta\,(\langle b_0,e\rangle + 1)\bigr|.
\]

\section{Joint $(k,\tau)$ selection with $\theta = \mu = 1/k^2$}

Setting $\theta = \mu = 1/k^2$ ties the centering weight to the barrier
parameter.  The direction becomes
\[
  d = d_0 + k\tau\,d_1^{(0)} + \frac{1}{k}\,d_1^{(\theta)},
\]
which depends on two free parameters $(k, \tau)$.

\subsection{Inner-product basis}

All norms and inner products of $d$ can be expressed in terms of
six scalars computed once from the decomposition:
\begin{alignat*}{3}
  &a = \|d_0\|^2, &\qquad
  &b = \langle d_0,\, d_1^{(0)}\rangle, &\qquad
  &c = \langle d_0,\, d_1^{(\theta)}\rangle, \\
  &p = \|d_1^{(0)}\|^2, &\qquad
  &q = \langle d_1^{(0)},\, d_1^{(\theta)}\rangle, &\qquad
  &r = \|d_1^{(\theta)}\|^2.
\end{alignat*}

\subsection{Optimal $\tau$ at fixed $k$}

The squared norm
\[
  \|d\|^2 = a + 2k\tau b + \frac{2c}{k}
            + k^2\tau^2 p + 2\tau q + \frac{r}{k^2}
\]
is quadratic in~$\tau$.  Setting $\partial\|d\|^2/\partial\tau = 0$:
\[
  \boxed{\tau^*(k) = -\frac{k\,b + q}{k^2\,p}.}
\]

\subsection{Reduced objective in $k$ alone}

Substituting $\tau^*(k)$ back:
\[
  \|d\|^2\bigl(k,\,\tau^*(k)\bigr)
  = a + \frac{2c}{k} + \frac{r}{k^2}
    - \frac{(k\,b + q)^2}{k^2\,p}.
\]
This is a rational function of~$k$ involving only the six precomputed
scalars.  It is minimized by a one-dimensional search over $k > 0$.

\subsection{Algorithm (per factorization)}

\begin{enumerate}
  \item Compute the decomposition $(d_0,\, d_1^{(0)},\, d_1^{(\theta)})$:
        one factorization and three back-solves.
  \item Compute the six inner products $a, b, c, p, q, r$.
  \item Line-search over $k > 0$ to minimize $\|d\|^2(k, \tau^*(k))$.
  \item Set $\tau = \tau^*(k)$, $\theta = 1/k^2$, $\mu = 1/k^2$.
  \item Evaluate $d = d_0 + k\tau\,d_1^{(0)} + (1/k)\,d_1^{(\theta)}$.
  \item Take the geodesic step $W \leftarrow W \cdot \exp(\alpha\,d)$
        with $\alpha = \min(1,\; 2/\|d\|_\infty^2)$.
\end{enumerate}

No additional Gram solves are required beyond the initial decomposition.
The entire parameter selection is determined by inner products of the
precomputed direction components.

\section{Derivation of the decomposition}

Starting from the single-parameter blend
$b_\theta = \theta\,e + (1-\theta)\,b_0$,
$c_\theta = \theta\,A^T e + (1-\theta)\,c$,
the combined RHS for $y_1$ is
\begin{align*}
  G\,y_1
  &= -\bigl(c_\theta + A^T P(W)\,b_\theta\bigr) \\
  &= -\bigl(c + A^T P(W)\,b_0\bigr)
     + \theta\,\bigl(c + A^T P(W)\,b_0 - A^T e - A^T P(W)\,e\bigr),
\end{align*}
giving $y_1 = y_1^{(0)} + \theta\,y_1^{(\theta)}$.
Expanding $d_1 = \PW(-b_\theta - A\,y_1)$ yields
$d_1 = d_1^{(0)} + \theta\,d_1^{(\theta)}$, which corresponds to
$\tau = 1 - \theta$.  Decoupling $\tau$ and $\theta$ generalizes this
to the full two-parameter family.

\end{document}
